\section{Mixed Dougall derivatives at a double residue zero}
\label{dg:section}

The centered resonance theorem already evaluates the minimal
subtraction at every nonpositive integer excess parameter. Its
residue-zero loci contain further identities: after restriction to
some loci the remainder vanishes term by term, while derivatives
transverse to those loci produce nonzero sums. We make this statement
precise, give a finite formula at every mixed order, and evaluate a
weighted square of trigamma differences for an arbitrary positive
shift. An additional spectral variable gives an all-order generating
function for generalized Stieltjes constants.

The underlying summation is the classical Rogers--Dougall
${}_5F_4(1)$ formula, DLMF 16.4.9~\cite{DLMFDougall}. The summation itself and the
centered subtraction theorem are antecedents, not novelty claims~\cite{RepoDougall,RepoTransport}.
The formulas below resolve the explicit mixed-parameter target
at $b=c=1$ in the incoming coordinate-transport report. No claim of
global literature priority is required for these deductions.

\subsection{Exact polynomial collapse on the integer grid}

Use the centered notation
\[
 \eta=a/2,\quad x_n=n+\eta,\quad d_0=1+a-b-c,
\]
and
\[
 r_n(u)=x_n\frac{\Gamma(n+a)\Gamma(n+b)\Gamma(n+c)
                     \Gamma(n+d_0-u)}
 {\Gamma(n+1)\Gamma(n+1+a-b)\Gamma(n+1+a-c)
                     \Gamma(n+b+c+u)}.
\]
Reciprocal Gamma factors are interpreted as entire functions, including
at their zeros. Put
\[
 (\alpha_1,\alpha_2,\alpha_3,\alpha_4)
       =(a/2,b-a/2,c-a/2,d_0-u-a/2),
\]
and define finite polynomial data by
\begin{align}
 h_j(u)&=-\frac{\sum_{i=1}^4B_{2j+1}(\alpha_i)}{j(2j+1)},\qquad j\geq1,
       \notag\\
 C_0(u)&=1,\qquad
 jC_j(u)=\sum_{k=1}^jkh_k(u)C_{j-k}(u).
 \label{dg:centered-recurrence}
\end{align}
Here $B_j(z)$ denotes the Bernoulli polynomial with $B_1(z)=z-1/2$.
These are exactly the centered asymptotic coefficients:
\[
 r_n(u)\sim x_n^{-1-2u}\sum_{j\geq0}C_j(u)x_n^{-2j},\qquad C_0=1.
\]
At $u=-N$, $N\geq1$, define the normally convergent bracket
\[
 F_{N,n}(a,b,c)=r_n(-N)-\sum_{j=0}^N C_j(-N)x_n^{2N-1-2j}.
\]

\begin{lemma}[Centered summation and residue transport: antecedent]
\label{dg:antecedent}
Suppose $a,b,c>0$ and $d_0>0$. For $K\geq1$ and
$-K<\Re u<d_0$, the normally convergent series
\begin{align}
 R_K(u)&=\sum_{n\geq0}\left\{r_n(u)-
       \sum_{j=0}^{K-1}C_j(u)x_n^{-1-2u-2j}\right\}\notag\\
 &=G(u)-\sum_{j=0}^{K-1}C_j(u)\zeta(1+2u+2j,a/2),\notag\\
 G(u)&=\frac{\Gamma(b)\Gamma(c)\Gamma(d_0-u)\Gamma(u)}
          {2\Gamma(d_0)\Gamma(b+u)\Gamma(c+u)}.
 \label{dg:antecedent-master}
\end{align}
The combined right side is holomorphic throughout this strip.
For $N\geq0$, let
\[
 \rho_N=\frac{(-1)^N}{N!}(d_0)_N(1-b)_N(1-c)_N.
\]
Then
\begin{align}
 C_N(-N)&=\rho_N,\notag\\
 R_{N+1}(-N)&=\rho_N\left\{\psi(a/2)+
 \frac{\psi(N+1)-\psi(b)-\psi(c)-\psi(d_0+N)}2\right\}.
 \label{dg:antecedent-closed}
\end{align}
All fixed spectral and parameter derivatives may be taken in the
single bracket in \eqref{dg:antecedent-master}.
\end{lemma}
\begin{proof}
At $x_n=n+a/2$, the denominator shifts are $1-\alpha_i$, so
\[
 r_n(u)=x_n\prod_{i=1}^4
                   \frac{\Gamma(x_n+\alpha_i)}{\Gamma(x_n+1-\alpha_i)}.
\]
The finite Stirling--Bernoulli expansion~\cite{DLMFGammaAsymptotic} and Bernoulli reflection
give, uniformly on compact parameter sets,
\[
 r_n(u)=x_n^{-1-2u}\left\{\sum_{j=0}^{K-1}C_j(u)x_n^{-2j}
                                     +O(n^{-2K})\right\}.
\]
Only even inverse powers occur; exponentiating the finite logarithmic
expansion gives \eqref{dg:centered-recurrence}. Thus the bracket is
summable throughout $-K<\Re u<d_0$. Cauchy's formula on a slightly
larger compact set proves the same for each derivative, with fixed
logarithmic powers allowed. In $0<\Re u<d_0$, the classical
Rogers--Dougall formula~\cite{DLMFDougall} sums $r_n(u)$ to $G(u)$:
the upper parameters are $a,1+a/2,b,c,d_0-u$, the lower parameters
are $a/2,1+a-b,1+a-c,b+c+u$, and the conventional excess is $2u$.
Separate summation of the counterterms gives
\eqref{dg:antecedent-master} there. Holomorphic continuation proves
the full-strip identity and removability of its apparent poles.

At $u=-N$ its Gamma residue is $\rho_N/2$. The only polar Hurwitz
counterterm has index $N$ and residue $C_N(-N)/2$. Equality of
residues gives $C_N(-N)=\rho_N$; since these are polynomial
coefficients, this equality extends through every residue zero.

For the constant term use
\[
 V=2\partial_a+\partial_b+\partial_c-\partial_u,
 \qquad W=\partial_a+\tfrac12\partial_b+\tfrac12\partial_c.
\]
One has $V\alpha_1=1$ and $V\alpha_i=0$ for $i=2,3,4$.
Consequently
\[
 Vh_k(u)=-\frac{B_{2k}(a/2)}k=2\zeta(1-2k,a/2),
 \qquad VC_j(u)=2\sum_{k=1}^j\zeta(1-2k,a/2)C_{j-k}(u).
\]
The second formula follows by differentiating the formal exponential
$\sum_jC_jy^j=\exp(\sum_kh_ky^k)$; every extracted coefficient is
finite. At $u=-N$ it becomes
\begin{equation}
 \frac12C_N'(-N)+\sum_{j=0}^{N-1}
 C_j(-N)\zeta(1-2N+2j,a/2)=W\rho_N.
 \label{dg:residue-transport}
\end{equation}
For generic $b,c$, the Gamma term has constant coefficient
\[
 \frac{\rho_N}{2}
 \{\psi(N+1)-\psi(d_0+N)-\psi(b-N)-\psi(c-N)\}.
\]
The resonant Hurwitz term has constant coefficient
$C_N'(-N)/2-\rho_N\psi(a/2)$. Use
\eqref{dg:residue-transport} to combine this derivative with all
lower counterterms. Since $Wd_0=0$,
\[
 W\rho_N=\frac{\rho_N}{2}
 \{\psi(b)-\psi(b-N)+\psi(c)-\psi(c-N)\}.
\]
The shifted digammas cancel and yield
\eqref{dg:antecedent-closed}. Its final arguments are positive, so
normal convergence and holomorphy extend it across all zeros of
$\rho_N$. These antecedent formulas reproduce the needed part of
the earlier resonance construction~\cite{RepoDougall,RepoTransport}.
\end{proof}

\begin{proposition}[Termwise collapse on every integer residue component]
\label{dg:polynomial}
Suppose $a,b,c>0$ and $d_0>0$. For $1\leq k\leq N$, at $c=k$,
\begin{align}
 r_n(-N)&=x_n
  \prod_{r=1}^{k-1}\bigl(x_n^2-(\eta-r)^2\bigr)
  \prod_{r=0}^{N-k-1}
        \bigl(x_n^2-(\eta+1-b+r)^2\bigr).
 \label{dg:grid-factor}
\end{align}
Consequently $F_{N,n}(a,b,k)=0$ for every $n$, and the same statement
holds with $b$ and $c$ interchanged. Locally along any intersection
$b=h$, $c=k$ with $h,k\in\{1,\ldots,N\}$, the holomorphic function
$F_{N,n}$ is divisible by $(b-h)(c-k)$. Its quotient and every fixed
derivative admit locally normally convergent sums. More globally, the
quotient by $(1-b)_N(1-c)_N$ has a holomorphic extension throughout
the stated parameter domain, and its sum is
\begin{align}
 &\sum_{n\geq0}\frac{F_{N,n}(a,b,c)}{(1-b)_N(1-c)_N}\notag\\
 &\quad=\frac{(-1)^N(d_0)_N}{N!}
 \left\{\psi(a/2)+
 \frac{\psi(N+1)-\psi(b)-\psi(c)-\psi(d_0+N)}2\right\}.
 \label{dg:divided-grid}
\end{align}
At a zero of either denominator factor, the summands in this formula
mean their holomorphic extensions just described.
\end{proposition}
\begin{proof}
The Gamma functional equation gives
\[
 r_n(-N)=x_n(n+1)_{k-1}(n+1+a-k)_{k-1}
                 (n+1+a-b)_{N-k}(n+b+k-N)_{N-k}.
\]
Pair the factors symmetric about $x_n$ to obtain
\eqref{dg:grid-factor}. This is an odd polynomial of degree
$2N-1$ in $x_n$. Uniqueness of asymptotic coefficients shows that its
centered coefficients are exactly $C_0(-N),\ldots,C_{N-1}(-N)$,
whereas $C_N(-N)=0$ and every subsequent coefficient is zero.
The identity includes the finitely many small indices where an
individual reciprocal Gamma factor vanishes.

For the divisibility statement, the holomorphic function vanishes
when $b=h$ and when $c=k$. Successive one-variable Taylor expansions
therefore factor $(b-h)(c-k)$. Cauchy's formula on a slightly larger
parameter polydisc transfers the $O(n^{-3})$ uniform remainder bound
to the divided function and to every prescribed derivative. This also
proves normal convergence after division; a quotient cannot be justified
merely by dividing two limiting sums at their common zero.
Every zero of $(1-b)_N$ or $(1-c)_N$ is simple. The other local
factors are units, so this establishes extension of the full divided
sum. Its value follows from \eqref{dg:antecedent-closed} off the
zero loci and then by holomorphy on them.
\end{proof}

The result is stronger than the vanishing of the residue polynomial
or of the summed remainder. In particular, every derivative tangent
to one of these components still has zero remainder termwise.
Mixed transverse derivatives are the first potentially nonzero ones.
For example, write
\[
 B_N(a,b,c)=\psi(a/2)+
 \frac{\psi(N+1)-\psi(b)-\psi(c)-\psi(d_0+N)}2.
\]
At $b=h$, $c=k$ with $1\leq h,k\leq N$,
\begin{align}
 \sum_{n\geq0}\partial_b\partial_cF_{N,n}(a,h,k)
 ={}&\frac{(-1)^{N+h+k}(d_0)_N}{N!}
 (h-1)!(N-h)!(k-1)!(N-k)!\notag\\
 &\times B_N(a,h,k),\qquad d_0=1+a-h-k.
 \label{dg:grid-first-mixed}
\end{align}
This uses $\left.\partial_b(1-b)_N\right|_{b=h}
=(-1)^h(h-1)!(N-h)!$ and keeps all finite-head terms.

\subsection{A two-parameter generating identity for all positive shifts}

Write $b=1+p$, $c=1+q$, $A=a-1$, and keep $u=-1$. Put
\begin{align}
 Q_n(a;p,q)
 &=\frac{\Gamma(n+a)\Gamma(n+1+p)\Gamma(n+1+q)
                         \Gamma(n+a-p-q)}
 {\Gamma(n+1)\Gamma(n+a-p)\Gamma(n+a-q)
                         \Gamma(n+1+p+q)},\label{dg:Q}\\
 B_a(p,q)&=\psi(a/2)+\frac{\psi(2)-\psi(1+p)-\psi(1+q)
                                      -\psi(a-p-q)}2.
 \label{dg:B}
\end{align}
Thus $Q_n(a;p,0)=Q_n(a;0,q)=1$, and $r_n(-1)=x_nQ_n(a;p,q)$.

\begin{theorem}[Double-zero generating function]
\label{dg:generator}
For every $a>0$ and sufficiently small complex $p,q$,
\begin{align}
 \sum_{n=0}^{\infty}\left\{
 x_n\bigl(Q_n(a;p,q)-1\bigr)
       +\frac{(A-p-q)pq}{x_n}\right\}
 =-(A-p-q)pq\,B_a(p,q).
 \label{dg:generator-eq}
\end{align}
The series is locally normally convergent in $(a,p,q)$ in a
neighborhood of each point $(a,0,0)$ with $\Re a>0$. It may be
differentiated at every fixed mixed order. Both sides are divisible
by $pq$, and the divided identity also holds normally at $p=q=0$.
\end{theorem}
\begin{proof}
For $a>1$ and small $p,q$, Lemma~\ref{dg:antecedent} applies with
$N=1$. Its residue coefficient is
\[
 C_1(-1)=-(a-1-p-q)pq,
\]
and the closed resonance value is exactly the right side of
\eqref{dg:generator-eq}. The centered expansion supplies
\[
 Q_n(a;p,q)=1-\frac{(A-p-q)pq}{x_n^2}+O(n^{-4})
\]
locally uniformly, so the displayed bracket is $O(n^{-3})$.
All Gamma arguments in \eqref{dg:Q} lie in a pole-free
neighborhood when $\Re a>0$ and $p,q$ are sufficiently small; the
same centered expansion is uniform on compact subsets there.
Thus the sum and the right side are holomorphic across $a=1$ and
throughout this connected positive-half-plane germ. The identity
theorem extends their equality from $a>1$ to all $a>0$.
For any prescribed compact $a$-set one can choose a smaller common
polydisc in $p,q$ and apply Cauchy's coefficient formula to the
uniform remainder. This proves the derivative and division claims.
\end{proof}

The restriction $d_0>0$ was needed only in the original region from
which continuation was proved. After the exact cancellation above,
$a=1$ is a regular point, not an excluded parameter or a formal limit.

\subsection{A finite Bell formula at every mixed order}

For $k\geq2$, define
\[
 \Lambda_k(n;a)=(-1)^k\psi^{(k-1)}(n+a)
                                -\psi^{(k-1)}(n+1).
\]
Taylor expansion of the four Gamma ratios gives the exact germ
\begin{equation}
 \log Q_n(a;p,q)=\sum_{r,s\geq1}
       \Lambda_{r+s}(n;a)\frac{p^rq^s}{r!s!}.
 \label{dg:logQ}
\end{equation}
Define a finite bivariate Bell polynomial by
\[
 \mathcal B_{rs}(n;a)=r!s![p^rq^s]
 \exp\left(\sum_{i,j\geq1}
       \Lambda_{i+j}(n;a)\frac{p^iq^j}{i!j!}\right).
\]
Only $i\leq r$, $j\leq s$, and powers at most $\min(r,s)$
of the exponent are needed. In particular,
\begin{align*}
 \mathcal B_{11}&=\Lambda_2,&
 \mathcal B_{21}&=\Lambda_3,&
 \mathcal B_{31}&=\Lambda_4,\\
 \mathcal B_{22}&=\Lambda_4+2\Lambda_2^2,&
 \mathcal B_{32}&=\Lambda_5+6\Lambda_2\Lambda_3,\\
 \mathcal B_{33}&=\Lambda_6+9\Lambda_2\Lambda_4
                            +9\Lambda_3^2+6\Lambda_2^3.
\end{align*}

Set $B_{rs}=\partial_p^r\partial_q^sB_a(0,0)$, with
\[
 B_{00}=\psi(a/2)+\frac{1+\gamma-\psi(a)}2,
\]
and, for $r+s=k\geq1$,
\begin{equation}
 B_{rs}=-\frac12\left\{
 \mathbf1_{s=0}\psi^{(r)}(1)+
 \mathbf1_{r=0}\psi^{(s)}(1)+(-1)^k\psi^{(k)}(a)\right\}.
 \label{dg:B-jets}
\end{equation}
Terms carrying a negative subscript below are omitted.

\begin{corollary}[Complete mixed parameter table]
\label{dg:all-parameter}
For $r,s\geq1$, the single bracketed series
\begin{align}
 S_{rs}(a)=\sum_{n\geq0}\left\{
 x_n\mathcal B_{rs}(n;a)
 +\frac{A\mathbf1_{r=s=1}
  -2\mathbf1_{(r,s)=(2,1)}-2\mathbf1_{(r,s)=(1,2)}}{x_n}
 \right\}
 \label{dg:S}
\end{align}
converges absolutely, and
\begin{equation}
 S_{rs}(a)=-rs\left\{
 A B_{r-1,s-1}-(r-1)B_{r-2,s-1}-(s-1)B_{r-1,s-2}
 \right\}.
 \label{dg:all-parameter-eq}
\end{equation}
For every $r+s\geq4$, no counterterm is needed: the original
polygamma polynomial $x_n\mathcal B_{rs}(n;a)$ is summable by itself.
\end{corollary}
\begin{proof}
Differentiate \eqref{dg:generator-eq} and use
\eqref{dg:logQ}. The counterterm is a polynomial of total
degree three in $p,q$, which explains its exact disappearance at
total derivative order four. Leibniz's rule gives
\eqref{dg:all-parameter-eq}; normal convergence was proved
before differentiation, including the small indices.
\end{proof}

For example, at $a=2$ the first rows are
\[
\begin{array}{c|c|c}
 (r,s)&\mathcal B_{rs}&S_{rs}(2)\\\hline
 (1,1)&\Lambda_2&0\\
 (2,1)&\Lambda_3&1\\
 (3,1)&\Lambda_4&-6\zeta(3)\\
 (2,2)&\Lambda_4+2\Lambda_2^2&-4\zeta(3)\\
 (3,2)&\Lambda_5+6\Lambda_2\Lambda_3&24\zeta(3)-18\zeta(4)\\
 (3,3)&\Lambda_6+9\Lambda_2\Lambda_4+9\Lambda_3^2+6\Lambda_2^3
                                    &108(\zeta(4)-\zeta(5)).
\end{array}
\]
The first two rows use the counterterms in \eqref{dg:S}.
Every later row is an ordinary convergent sum without subtraction.

\subsection{A closed square of trigamma differences}

\begin{theorem}[Weighted trigamma square]
\label{dg:square}
For every real $a>0$,
\begin{align}
 &\sum_{n=0}^{\infty}\left(n+\frac a2\right)
       \bigl(\psi'(n+a)-\psi'(n+1)\bigr)^2\notag\\
 &\qquad=\frac{a-1}{4}\psi''(a)
     +\frac{\psi'(a)-\zeta(2)}2
     +\frac{3(a-1)}2\zeta(3).
 \label{dg:square-eq}
\end{align}
The series converges locally normally for $\Re a>0$, with an
ordinary algebraic square in that complex extension.
\end{theorem}
\begin{proof}
The $(2,2)$ and $(3,1)$ rows give respectively
\begin{align*}
 \sum_nx_n(\Lambda_4+2\Lambda_2^2)
   &=2A\psi''(a)+4\psi'(a)-4\zeta(2),\\
 \sum_nx_n\Lambda_4
   &=\frac{3A}{2}\{\psi''(1)+\psi''(a)\}
                                  +3\{\psi'(a)-\zeta(2)\}.
\end{align*}
Subtract and divide by two, using $\psi''(1)=-2\zeta(3)$.
Every series used in the subtraction is absolutely convergent.
Independently, the polygamma asymptotic or its convergent series
shows $\Lambda_2(n;a)=O(n^{-2})$ uniformly on compact subsets of
$\Re a>0$. Thus the weighted square is $O(n^{-3})$; this also
directly verifies the domain and regularity at $a=1$.
\end{proof}

At $a=3/2$, the functional equation and duplication formula~\cite{DLMFPolygamma} give
\[
 \psi'(3/2)=3\zeta(2)-4,\qquad
 \psi''(3/2)=-14\zeta(3)+16.
\]
The right side of \eqref{dg:square-eq} becomes
\[
 \frac18(-14\zeta(3)+16)
 +\frac12(2\zeta(2)-4)+\frac34\zeta(3)
 =\zeta(2)-\zeta(3).
\]
For $H_n^{(r)}=\sum_{k=1}^nk^{-r}$, $H_0^{(r)}=0$,
\[
 \psi'(n+1)-\psi'(n+3/2)
 =2\{2H_{2n+1}^{(2)}-H_n^{(2)}-\zeta(2)\}.
\]
Therefore
\begin{equation}
 \boxed{\sum_{n=0}^{\infty}(4n+3)
 \{2H_{2n+1}^{(2)}-H_n^{(2)}-\zeta(2)\}^{2}
 =\zeta(2)-\zeta(3).}
 \label{dg:half-harmonic-square}
\end{equation}
No regularized interpretation is involved in this positive series.

The continuation to $a=1/2$ gives a companion identity,
\begin{equation}
 \sum_{n=0}^{\infty}(4n+1)
 \{\zeta(2)-2H_{2n}^{(2)}+H_n^{(2)}\}^{2}
 =\zeta(2)+\zeta(3).
 \label{dg:other-half-harmonic-square}
\end{equation}
Here $\psi'(1/2)=3\zeta(2)$ and
$\psi''(1/2)=-14\zeta(3)$. At $a=2$, the square formula reduces to
$\sum_{k\geq1}k^{-3}=\zeta(3)$, a useful elementary normalization.

The $(3,3)$ row also supplies a higher harmonic example. At $a=2$,
writing $k=n+1$,
\[
 \Lambda_2=-k^{-2},\quad \Lambda_4=-6k^{-4},\quad
 \Lambda_6=-120k^{-6},\quad
 \Lambda_3=2\{2(\zeta(3)-H_{k-1}^{(3)})-k^{-3}\}.
\]
Substitution gives
\begin{equation}
 \boxed{\sum_{k=1}^{\infty}k
 \{2(\zeta(3)-H_{k-1}^{(3)})-k^{-3}\}^2
 =3\zeta(4)-\zeta(5).}
 \label{dg:cubic-tail-square}
\end{equation}

\subsection{Every resonant polynomial moment from the same germ}

For $N\geq1$, $s=p+q$, set
\[
 P_{N-1}(x;s)=\prod_{r=0}^{N-2}\{x^2-(a/2-s+r)^2\},\qquad P_0=1.
\]
The exact Gamma functional equation gives
\begin{equation}
 r_n(-N)=x_nP_{N-1}(x_n;p+q)Q_n(a;p,q).
 \label{dg:N-transport}
\end{equation}
It retains every finite-head zero as a factor before derivatives are
taken. This identity and the two-variable Bell germ generate the
summands at every resonance without singular logarithmic derivatives.

More explicitly, write
\[
 P_{N-1}(x;s)=\sum_{k=0}^{N-1}P_k(s)x^{2N-2-2k},\qquad
 Q_n(a;p,q)\sim\sum_{j\geq0}D_j(p,q)x_n^{-2j}.
\]
Here $D_0=1$ and every $D_j$, $j\geq1$, is divisible by $pq$.
The counterterms are the finite convolution
\begin{equation}
 C_j(-N)=\sum_{k=0}^{\min(j,N-1)}P_k(p+q)D_{j-k}(p,q).
 \label{dg:N-coefficients}
\end{equation}
Equation~\eqref{dg:antecedent-closed} becomes
\begin{align}
 \sum_{n\geq0}F_{N,n}(a,1+p,1+q)
 &=\frac{(-1)^N}{N!}(a-1-p-q)_N(-p)_N(-q)_N\notag\\
 &\quad\times B_N(a,1+p,1+q).
 \label{dg:N-all}
\end{align}
This holds as a normally convergent germ for every $a>0$, by the
same continuation argument as Theorem~\ref{dg:generator}.
Taking any prescribed $[p^rq^s]$ is a finite exact algorithm involving
only the coefficients in \eqref{dg:N-coefficients}, polygamma
values, and finite generalized harmonic numbers from the Pochhammer
factors. Formula \eqref{dg:N-all} is a representation theorem;
it does not assert that all possible products of polygamma functions
can be separated individually by these mixed derivatives.

\subsection{A mixed spectral germ containing every Stieltjes order}

We now keep one derivative in each parameter direction and permit
the spectral parameter to move: $u=-1+t$. For $a>0$, let
\begin{align}
 R_n(a;t)&=x_n\frac{\Gamma(n+1)\Gamma(n+a-t)}
                       {\Gamma(n+a)\Gamma(n+1+t)},\notag\\
 L_n(a;t)&=\psi(n+1)-\psi(n+a-t)+\psi(n+a)-\psi(n+1+t),\notag\\
 K_n(a;t)&=\psi'(n+a-t)-\psi'(n+1+t).
 \label{dg:spectral-summands}
\end{align}
At $p=q=0$, these satisfy
\[
 \partial_p\partial_q r_n(-1+t)
                 =R_n(a;t)\{L_n(a;t)^2+K_n(a;t)\}.
\]
The identity uses the unspecialized summand before derivatives.

To state the endpoint without a division by $a-1$, define
\begin{align}
 F_a(t)&=\frac{\Gamma(a-t)}{\Gamma(a)\Gamma(1+t)(1-t)},\notag\\
 V_a(t)&=\psi(1)-\psi(a-t)+\psi(a)-\psi(1+t),\notag\\
 W_a(t)&=\psi'(a-t)-\psi'(a),\notag\\
 D_a(t)&=F_a(t)\left\{2t(1+tV_a(t))
 -(a-1)\bigl[(1+tV_a(t))^2+t^2W_a(t)\bigr]\right\}.
 \label{dg:D}
\end{align}
All of these are holomorphic near $t=0$ for every $a>0$, including
$a=1$, and $D_a(0)=-(a-1)$.

\begin{theorem}[Complete mixed spectral--Stieltjes germ]
\label{dg:spectral}
For $a>0$ and sufficiently small complex $t$, the series
\begin{equation}
 M_a(t)=\sum_{n=0}^{\infty}\left\{
 R_n(a;t)\bigl(L_n(a;t)^2+K_n(a;t)\bigr)
 +(a-1-2t)x_n^{-1-2t}\right\}
 \label{dg:M}
\end{equation}
converges locally normally and equals
\begin{equation}
 M_a(t)=\frac{D_a(t)}{2t}
             +(a-1-2t)\zeta(1+2t,a/2).
 \label{dg:M-eq}
\end{equation}
The apparent singularity at $t=0$ is removable. Define $\gamma_m(v)$
by the Laurent expansion at the simple Hurwitz pole~\cite{DLMFHurwitz}:
\[
 \zeta(1+z,v)=z^{-1}+\sum_{m\geq0}\frac{(-1)^m}{m!}\gamma_m(v)z^m.
\]
Then every fixed derivative has the explicit expression
\begin{equation}
 M_a^{(m)}(0)=\frac{D_a^{(m+1)}(0)}{2(m+1)}
 -\mathbf1_{m=0}
 +(-2)^m\{(a-1)\gamma_m(a/2)+m\gamma_{m-1}(a/2)\}.
 \label{dg:all-Stieltjes}
\end{equation}
The $m\gamma_{m-1}$ term is omitted when $m=0$. Every derivative
on the left can be taken in the single convergent bracket
\eqref{dg:M}.
\end{theorem}
\begin{proof}
For $a>1$, differentiate \eqref{dg:antecedent-master} with $K=2$
once in $p$ and once in $q$. The lower counterterm $C_0=1$ is
parameter independent. Direct differentiation of the cubic Bernoulli
coefficient gives
\[
 \left.\partial_p\partial_qC_1(-1+t)\right|_{p=q=0}=-(a-1)+2t.
\]
The Gamma term is $E_1(t;p,q)/(2t)$, where
\[
 E_1(t;p,q)=
 \frac{\Gamma(1+p)\Gamma(1+q)\Gamma(a-p-q-t)\Gamma(1+t)}
 {\Gamma(a-1-p-q)\Gamma(1+p+t)\Gamma(1+q+t)}
                  \frac{(p+t)(q+t)}{t-1}.
\]
The factors $(p+t)(q+t)$ must be retained: setting $p=q=0$
before differentiating would erase the mixed derivative.
Twice applying the product rule and using
\[
 \Gamma(a)=(a-1)\Gamma(a-1),\quad
 \psi(a-1)=\psi(a)-(a-1)^{-1},\quad
 \psi'(a-1)=\psi'(a)+(a-1)^{-2}
\]
shows that $\partial_p\partial_qE_1(t;0,0)=D_a(t)$.
The resulting expression \eqref{dg:D} contains no singular
factor at $a=1$. This proves \eqref{dg:M-eq} for $a>1$.
Uniform Stirling remainders on a slightly enlarged parameter
neighborhood give a bound
$O(n^{-3-2\Re t}(1+\log^j n))$ for every fixed $t$-derivative of
the bracket. The displayed functions are holomorphic for $\Re a>0$
near $t=0$, so continuation gives all $a>0$ as before.

Finally subtract the pole explicitly:
\[
 M_a(t)=\frac{D_a(t)+(a-1)-2t}{2t}
 +(a-1-2t)\sum_{h\geq0}\frac{(-2)^h}{h!}\gamma_h(a/2)t^h.
\]
Taylor coefficient extraction yields
\eqref{dg:all-Stieltjes}, including the constant $-1$ at
$m=0$. The formula is for derivatives rather than divided Taylor
coefficients; its factor $1/(m+1)$ is therefore essential.
\end{proof}

For an explicit first spectral derivative, write
\[
 \ell=1+\gamma-\psi(a),\quad
 d=\psi'(a)-\zeta(2),\quad
 \mu_n=-\psi(n+a)-\psi(n+1).
\]
Then \eqref{dg:all-Stieltjes} gives the ordinary identity
\begin{align}
 &\sum_{n=0}^{\infty}\left\{
 x_n\bigl[\Lambda_3(n;a)+\mu_n\Lambda_2(n;a)\bigr]
       -\frac{2+2(a-1)\log x_n}{x_n}\right\}\notag\\
 &\quad=2\psi(a/2)+\ell
       -\frac{a-1}{4}\bigl(\ell^2+5d+1\bigr)
       -2(a-1)\gamma_1(a/2).
 \label{dg:first-spectral}
\end{align}
Indeed the coefficient of $t^2$ in $D_a(t)$ is
$2\ell-(a-1)(\ell^2+5d+1)/2$.
At $a=2$, the right side reduces to
$1-\gamma^2-2\gamma_1$. Writing $k=n+1$, this is equivalently
\begin{align}
 \sum_{k=1}^{\infty}\left\{
 4k(\zeta(3)-H_{k-1}^{(3)})-\frac1{k^2}-\frac2k
 +\frac{2(H_{k-1}-\gamma-\log k)}k\right\}
 =1-\gamma^2-2\gamma_1.
 \label{dg:harmonic-Stieltjes}
\end{align}
The divergent-looking pieces in this last formula are one bracket;
the proof supplies absolute convergence of their combination.

\subsection{Further questions suggested by the mixed normal form}

The arbitrary-order formulas above solve the requested first
two-parameter table at $b=c=1$ and provide its spectral extension.
Several sharper questions remain. First, at a fixed homogeneous
order, which individual polygamma-product moments are recoverable
from the finitely many mixed Bell combinations? Their number need
not match the number of all possible products, so the all-order
generating theorem alone does not prove complete separation.
Second, the polynomial formula on $b=h$, $c=k$ supplies starting
points at every residue-grid intersection; compact explicit tables
away from $(1,1)$ should retain their finite-head zeros and avoid
singular logarithmic derivative notation. Third, at rational $a$,
the Stieltjes germ can be expressed in finite residue-class or
Dirichlet-character coordinates; it is natural to determine which
sectors cancel in particular mixed orders. Finally, primitive or
generating-variable versions of the mixed identities can be tested
for reductions to colored polylogarithms, with the function class
and endpoint normalization stated before any claimed reduction.
